\documentclass{article}
\usepackage{mathtools} 
\usepackage{fontspec}
\usepackage[UTF8]{ctex}
\usepackage{amsthm}
\usepackage{mdframed}
\usepackage{xcolor}
\usepackage{amssymb}
\usepackage{amsmath}


% 定义新的带灰色背景的说明环境 zremark
\newmdtheoremenv[
  backgroundcolor=gray!10,
  % 边框与背景一致，边框线会消失
  linecolor=gray!10
]{zremark}{说明}

% 通用矩阵命令: \flexmatrix{矩阵名}{元素符号}{行数}{列数}
\newcommand{\flexmatrix}[4]{
  \[
  #1 = \begin{pmatrix}
    #2_{11}     & #2_{12}     & \cdots & #2_{1#4}   \\
    #2_{21}     & #2_{22}     & \cdots & #2_{2#4}   \\
    \vdots      & \vdots      & \ddots & \vdots     \\
    #2_{#31}    & #2_{#32}    & \cdots & #2_{#3#4}
  \end{pmatrix}
  \]
}

% 简化版命令（默认矩阵名为A，元素符号为a）: \quickmatrix{行数}{列数}
\newcommand{\quickmatrix}[2]{\flexmatrix{A}{a}{#1}{#2}}

\begin{document}
\title{习题 12.1}
\author{张志聪}
\maketitle

\section*{1}

\begin{itemize}
  \item (1)

        我们有
        \begin{align*}
          S_n & = \frac{1}{5} (1 - \frac{1}{6}) + \frac{1}{5} (\frac{1}{6} - \frac{1}{11}) + \cdots + \frac{1}{5n - 4} (\frac{1}{6} - \frac{1}{5n + 1}) \\
              & = \frac{1}{5} \left( 1 - \frac{1}{5n + 1} \right)
        \end{align*}
        所以
        \begin{align*}
          \lim\limits_{n \to \infty} S_n & = \lim\limits_{n \to \infty} \frac{1}{5} \left( 1 - \frac{1}{5n + 1} \right) \\
                                         & = \frac{1}{5}
        \end{align*}

  \item (2)

        我们有
        \begin{align*}
          S_n & = (\frac{1}{2} + \cdots + \frac{1}{2^n}) + (\frac{1}{3} + \cdots + \frac{1}{3^n})                                   \\
              & = \frac{\frac{1}{2}(1- (\frac{1}{2})^n)}{1 - \frac{1}{2}} + \frac{\frac{1}{3}(1- (\frac{1}{3})^n)}{1 - \frac{1}{3}} \\
              & = 1- (\frac{1}{2})^n + \frac{1}{2}(1 - (\frac{1}{3})^n)
        \end{align*}
        注意：考虑部分和时，$S_n$是有限项求和，可以交换求和顺序，但级数本身在一般情况下是不可以的。

        所以
        \begin{align*}
          \lim\limits_{n \to \infty} S_n & = \lim\limits_{n \to \infty} 1- (\frac{1}{2})^n + \frac{1}{2}(1 - (\frac{1}{3})^n) \\
                                         & = \frac{3}{2}
        \end{align*}

  \item (3)

        方法一：

        设
        \begin{align*}
          \frac{A}{n} + \frac{B}{n + 1} + \frac{C}{n + 2} = \frac{1}{n (n + 1)(n + 2)}
        \end{align*}
        解出$A,B,C$的值为
        \begin{equation*}
          \begin{cases}
            A = \frac{1}{2} \\
            B = -1          \\
            C = \frac{1}{2}
          \end{cases}
        \end{equation*}
        于是
        \begin{align*}
          \frac{\frac{1}{2}}{n} - \frac{1}{n + 1} + \frac{\frac{1}{2}}{n + 2}
           & = \frac{\frac{1}{2}}{n} - \frac{\frac{1}{2}}{n + 1} - \frac{\frac{1}{2}}{n + 1} + \frac{\frac{1}{2}}{n + 2} \\
           & = \frac{1}{2}(\frac{1}{n} - \frac{1}{n + 1}) + \frac{1}{2}(\frac{1}{n + 2} - \frac{1}{n + 1})
        \end{align*}

        方法二：

        这个方法，计算量较小。

        我们有
        \begin{align*}
          \frac{1}{n (n + 1)(n + 2)}
           & = \frac{1}{n}\left(\frac{1}{(n + 1)(n + 2)}\right)                                               \\
           & = \frac{1}{n} \left( \frac{1}{n+1} - \frac{1}{n+2} \right)                                       \\
           & = \frac{1}{n(n+1)} - \frac{1}{n(n+2)}                                                            \\
           & = \left(\frac{1}{n} - \frac{1}{n+1}\right) - \frac{1}{2}\left(\frac{1}{n} - \frac{1}{n+2}\right)
        \end{align*}
        于是，我们有
        \begin{align*}
          S_n & = \left[\left(1 - \frac{1}{2}\right) + \cdots + \left(\frac{1}{n} - \frac{1}{n+1}\right)\right]                                                                                             \\
              & -\frac{1}{2}\left[\left(1 - \frac{1}{3}\right) + \left(\frac{1}{2} - \frac{1}{4}\right) + \left(\frac{1}{3} - \frac{1}{5}\right) + \cdots + \left(\frac{1}{n} - \frac{1}{n+2}\right)\right] \\
              & = 1 - \frac{1}{n+1} -\frac{1}{2}\left[\left(1 + \frac{1}{2} - \frac{1}{n+2}\right) \right]
        \end{align*}
        所以
        \begin{align*}
          \lim\limits_{n \to \infty} S_n & = 1 - \frac{1}{n+1} -\frac{1}{2}\left[\left(1 + \frac{1}{2} - \frac{1}{n+2}\right) \right] \\
                                         & = \frac{1}{4}
        \end{align*}

  \item (4)

        提示：
        \begin{align*}
          U_n = (\sqrt{n+2} - \sqrt{n+1}) + (\sqrt{n} - \sqrt{n+1})
        \end{align*}


  \item (5)

        我们有
        \begin{align*}
          S_n = \frac{1}{2} + \frac{3}{2^2} + \cdots + \frac{2n - 1}{2^n}
        \end{align*}
        可见，$S_n$的项是等差乘等比。

        于是，按照常用的处理办法：
        \begin{align*}
          \frac{1}{2}S_n & = \frac{1}{2^2} + \frac{3}{2^3} + \cdots + \frac{2n - 3}{2^n} + \frac{2n - 1}{2^{n+1}}
        \end{align*}
        两式相加，我们有
        \begin{align*}
          \frac{1}{2}S_n & = \frac{1}{2} + \frac{2}{2^2} + \frac{2}{2^3} + \cdots + \frac{2}{2^n} - \frac{2n - 1}{2^{n+1}}       \\
                         & = \frac{1}{2} + \frac{1}{2} + \frac{1}{2^2} + \cdots + \frac{1}{2^{n-1}} - \frac{2n - 1}{2^{n+1}}     \\
                         & = \frac{1}{2} + \frac{\frac{1}{2}(1 - (\frac{1}{2})^{n-1})}{1 - \frac{1}{2}} - \frac{2n - 1}{2^{n+1}} \\
                         & = \frac{3}{2} - (\frac{1}{2})^{n-1} - \frac{2n - 1}{2^{n+1}}
        \end{align*}
        所以
        \begin{align*}
          S_n = 3 - 2(\frac{1}{2})^{n-1} - \frac{2n - 1}{2^{n}}
        \end{align*}
        于是
        \begin{align*}
          \lim\limits_{n \to \infty} S_n & = \lim\limits_{n \to \infty} 3 - 2(\frac{1}{2})^{n-1} - \frac{2n - 1}{2^{n}} \\
                                         & = 3
        \end{align*}
\end{itemize}

\section*{3}

 (i) 不一定；

(i.1)比如
\begin{align*}
  \sum u_n & = \sum -1 \\
  \sum v_n & = \sum 1
\end{align*}
令
\begin{align*}
  \sum w_n = \sum (u_n + v_n) = \sum 0 = 0
\end{align*}
收敛；

(i.2)比如
\begin{align*}
  \sum u_n & = \sum 1 = \sum v_n
\end{align*}
令
\begin{align*}
  \sum w_n = \sum (u_n + v_n) = \sum 2
\end{align*}
发散。

(ii) 发散；

设
\begin{align*}
  \sum u_n, \sum v_n
\end{align*}
都是发散的非负数级数。

方法一：

可得$\sum u_n, \sum v_n$的部分和
\begin{align*}
  A_n = u_1 + u_2 + \cdots + u_n \\
  B_n = v_1 + v_2 + \cdots + v_n
\end{align*}
序列${A_n}, \{B_n\}$都是单调递增的，由于它们都是发散的，
所以它们无上界，有
\begin{align*}
  \sum (u_n + v_n)
\end{align*}
所以其部分和
\begin{align*}
  S_n & = (u_1 + v_1) + (u_2 + v_2) + \cdots + (u_n + v_n) \\
      & = A_n + B_n
\end{align*}
于是可知$\{S_n\}$是单调递增的，且没有上界，所以发散。

方法二：

由级数收敛的柯西准则，证明级数发散。

因为$\sum u_n$都是发散的，
所以存在$\epsilon_0 > 0$，对任意$N > 0$，总存在$m_0 > N, p_0 > 0$，有
\begin{align*}
  |u_{m_0 + 1} + u_{m_0 + 2} + \cdots + u_{m_0 + p_0}| > \epsilon_0
\end{align*}

于是，我们有
\begin{align*}
   & |(u_{m_0 + 1} + v_{m_0 + 1}) + (u_{m_0 + 2} + v_{m_0 + 2}) + \cdots + (u_{m_0 + p_0} + v_{m_0 + p_0})|        \\
   & = |u_{m_0 + 1} + u{m_0 + 2} + \cdots + u_{m_0 + p_0} + v_{m_0 + 1} + v_{m_0 + 2} + \cdots + v_{m_0 + p_0}|    \\
   & = |u_{m_0 + 1} + u{m_0 + 2} + \cdots + u_{m_0 + p_0} | + |v_{m_0 + 1} + v_{m_0 + 2} + \cdots + v_{m_0 + p_0}| \\
   & \geq |u_{m_0 + 1} + u_{m_0 + 2} + \cdots + u_{m_0 + p_0}| > \epsilon_0
\end{align*}
注意：以上用到了$u_n, v_n (n = 1, 2, \cdots)$都是非负数。

所以，$\sum (u_n + v_n)$发散。

\section*{4}

我们有
\begin{align*}
  S_n & = (a_1 - a_2) + (a_2 - a_3) + \cdots + (a_n - a_{n + 1}) \\
      & = a_1 - a_{n + 1}
\end{align*}
所以
\begin{align*}
  \lim\limits_{n \to \infty} S_n & = \lim\limits_{n \to \infty} a_1 - a_{n + 1} \\
                                 & = a_1 - a
\end{align*}

\section*{6}

\begin{itemize}
  \item (1)

        设数列${a_n}$为
        \begin{align*}
          a_n = \frac{1}{a + n - 1}
        \end{align*}
        显然，$\{a_n\}$收敛于$0$。\\
        又
        \begin{align*}
          \sum \limits_{n = 1}^\infty \frac{1}{(a + n - 1)(a + n)}
           & =  \sum \limits_{n = 1}^\infty \frac{1}{a + n - 1} - \frac{1}{a + n} \\
           & =  \sum \limits_{n = 1}^\infty a_n - a_{n + 1}
        \end{align*}
        由习题4可知，$\sum \limits_{n = 1}^\infty a_n - a_{n + 1} = a_1 - 0 = \frac{1}{a}$

  \item (2)

        我们有
        \begin{align*}
          (-1)^{n + 1}\frac{2n + 1}{n(n + 1)} & = (-1)^{n + 1}\frac{1}{n} + (-1)^{n + 1}\frac{1}{n + 1}           \\
                                              & = -(-1)^n \frac{1}{n} + (-1)^{n + 1}\frac{1}{n + 1}               \\
                                              & = - \left[(-1)^n \frac{1}{n} - (-1)^{n + 1}\frac{1}{n + 1}\right]
        \end{align*}
        设数列${a_n}$为
        \begin{align*}
          a_n = (-1)^n \frac{1}{n}
        \end{align*}
        显然，$\{a_n\}$收敛于$0$。\\
        又
        \begin{align*}
          \sum \limits_{n = 1}^\infty (-1)^{n + 1}\frac{2n + 1}{n(n + 1)}
           & =  \sum \limits_{n = 1}^\infty - (a_n - a_{n + 1}) \\
           & = - \sum \limits_{n = 1}^\infty (a_n - a_{n + 1})
        \end{align*}
        由习题4可知，$- \sum \limits_{n = 1}^\infty a_n - a_{n + 1} = -(a_1 - 0) = -(-1) = 1$.

  \item (3)

        我们有
        \begin{align*}
          \frac{2n + 1}{(n^2 + 1)[(n+1)^2 + 1]}
           & = \frac{1}{n^2 + 1} - \frac{1}{(n+1)^2 + 1}
        \end{align*}
        设数列${a_n}$为
        \begin{align*}
          a_n = \frac{1}{n^2 + 1}
        \end{align*}
        显然，$\{a_n\}$收敛于$0$。\\
        又
        \begin{align*}
          \sum \limits_{n = 1}^\infty \frac{2n + 1}{(n^2 + 1)[(n+1)^2 + 1]}
           & = \sum \limits_{n = 1}^\infty (a_n - a_{n + 1})
        \end{align*}
        由习题4可知，$\sum \limits_{n = 1}^\infty a_n - a_{n + 1} = a_1 - 0 = \frac{1}{2}$.
\end{itemize}

\section*{7}

\begin{itemize}
  \item (1)

        收敛；

        \begin{align*}
          |u_{m + 1} + u_{m + 2} + \cdots + u_{m + p}|
           & = |\frac{sin 2^{m + 1}}{2^{m + 1}} + \frac{sin 2^{m + 2}}{2^{m + 2}} + \cdots + \frac{sin 2^{m + p}}{2^{m + p}}| \\
           & \leq \frac{1}{2^{m + 1}} + \frac{1}{2^{m + 2}} + \cdots + \frac{1}{2^{m + p}}                                    \\
           & = \frac{1}{2^{m}} \left(\frac{1}{2} + \frac{1}{2^2} + \cdots + \frac{1}{2^p}\right)
        \end{align*}
        由$\sum \limits_{n \to \infty} \frac{1}{2^n} = 1$可知，
        \begin{align*}
          \frac{1}{2^{m}} \left(\frac{1}{2} + \frac{1}{2^2} + \cdots + \frac{1}{2^p}\right)
           & \leq \frac{1}{2^{m}}
        \end{align*}
        由$\lim\limits_{n \to \infty} \frac{1}{2^n}$可知，对任意$\epsilon > 0$，
        都存在$N$满足柯西准则要求，所以级数收敛。

  \item (2)

        发散；

        \begin{align*}
          \lim\limits_{n \to \infty} u_n
           & = \lim\limits_{n \to \infty} (-1)^{n-1} \frac{n^2}{2n^2 + 1} \\
           & = (-1)^{n-1} \frac{1}{2}
        \end{align*}
        由级数收敛的柯西准则的推论可知，级数发散。

  \item (3)

        我们有
        \begin{align*}
          |u_{m + 1} + u_{m + 2} + \cdots + u_{m + p}|
           & = \left|\frac{(-1)^{m + 1}}{m+1} + \frac{(-1)^{m + 2}}{m+2} + \cdots + \frac{(-1)^{m + p}}{m+p}\right|              \\
           & = \left|(-1)^{m + 1} (\frac{1}{m + 1} - \frac{1}{m + 2} + \cdots + \frac{(-1)^{p-1}}{m + p})\right|                 \\
           & = \left|\frac{1}{m + 1} - \frac{1}{m + 2} + \cdots + \frac{(-1)^{p-1}}{m + p}\right|                                \\
           & = \left|\frac{1}{m + 1} - \left(\frac{1}{m + 2} - \frac{1}{m + 3} + \cdots + \frac{(-1)^{p-2}}{m + p}\right)\right|
        \end{align*}
        因为
        \begin{align*}
           & \frac{1}{m + 2} - \frac{1}{m + 3} + \cdots + \frac{(-1)^{p-2}}{m + p}                                           \\
           & = (\frac{1}{m + 2} - \frac{1}{m + 3}) + (\frac{1}{m + 4} - \frac{1}{m + 5}) + \cdots + \frac{(-1)^{p-2}}{m + p} \\
           & > 0
        \end{align*}
        所以
        \begin{align*}
          \left|\frac{1}{m + 1} - \left(\frac{1}{m + 2} - \frac{1}{m + 3} + \cdots + \frac{(-1)^{p-2}}{m + p}\right)\right|
          < \frac{1}{m + 1}
        \end{align*}
        于是对任意$\epsilon > 0$，取$N = \frac{1}{\epsilon} - 1$，当$m > N$时，我们有$\frac{1}{m + 1} < \epsilon$。

  \item (4)

        当$p = m$时，我们有
        \begin{align*}
           & |u_{m + 1} + u_{m + 2} + \cdots + u_{m + p}|                                                                                                    \\
           & = \left| \frac{1}{\sqrt{(m + 1) + (m + 1)^2}} + \frac{1}{\sqrt{(m + 2) + (m + 2)^2}} + \cdots + \frac{1}{\sqrt{(m + m) + (m + m)^2}} \right|    \\
           & \geq \left| \frac{1}{\sqrt{(m + m) + (m + m)^2}} + \frac{1}{\sqrt{(m + m) + (m + m)^2}} + \cdots + \frac{1}{\sqrt{(m + m) + (m + m)^2}} \right| \\
           & = \left| \frac{m}{\sqrt{(m + m) + (m + m)^2}} \right|                                                                                           \\
           & > \frac{1}{2}
        \end{align*}
        所以，级数发散。
\end{itemize}

\section*{9}

调和级数。

\section*{10}




\end{document}